\section{Computation: what exact inference costs and how to cut it}
\label{sec:compute}

The advantage of the structured estimator is in data and understanding, not in wall-clock, and the cost
deserves its own measurement. Grid BP is the bottleneck of the method: each message update is a product of an
$\ngrid\times\ngrid$ matrix with $\ngrid$-vectors, once per site and per level. We measure it on exclusive
nodes (the whole node allocated), report the thread count with every number, and then ask whether the grid
can be smaller or the messages represented differently.

\paragraph{Cost and its scaling.}
On the paper's $401$-point grid with eight threads, one posterior mean of a $32$-site sequence takes
$\cmpBpMsEightThreads$\,ms against $\cmpTfMsEightThreads$\,ms for the largest transformer, a factor
$\cmpRatioCpu$ on a CPU node and $\cmpRatioGpu$ on an H200 GPU (at batch $512$; the GPU time is flat in
the batch because BP is bound by kernel launches, so the ratio falls to $\cmpRatioGpuLarge$ at batch $2048$).
The cost is linear in the number of sequences and in the chain length (log-log slopes $\cmpSlopeBatch$ and
$\cmpSlopeLength$) and grows like $\ngrid^{\cmpSlopeGrid}$ rather than $\ngrid^2$, because elementwise work and
memory traffic dominate the matrix product at these sizes. Training is not cheaper than a network's: the
E-step is linear in the number of sequences and the M-step does not depend on it, but total EM time is
comparable with a network's, not orders of magnitude smaller (\cref{app:cost}). Memory is five message-sized
arrays, $\cmpMemMb$\,MB for a $512\times32$ batch at $401$ points, linear in batch, length and grid, and
$\cmpMemRatioMin$--$\cmpMemRatioMax$ times the largest network.

\paragraph{The grid can be coarser.}
The $401$-point grid was chosen for convergence safety and checked from above only (finer grids change nothing).
\Cref{fig:coarsegrid} checks it from below. With the true $\phi^4$ kernel the score is exact to
$\cmpCoarseOracleHundred$ at $101$ points on every cell, because the trapezoid rule on an analytic integrand
converges faster than any power of the spacing; and the \emph{learned} $\phi^4$ kernels lose nothing: at $101$
points the site energy's and the spline's risks equal their $401$-point values to three digits on two cells and
at two sample sizes, and at $51$ points they move by at most $\cmpCoarseLearnedFifty$, while the EM fit is
$\cmpCoarseFitSpeedup$ times faster. On one thread this makes BP $\cmpCoarseSpeedupHundred$ times faster at $101$
points and $\cmpCoarseSpeedupOneFifty$ times at $151$ than at $401$. The linear chain's kernels are harder: with
Laplace noise the transition is not smooth and, with many sequences, the eight-component mixture learns components
narrower than a coarse spacing, so its risk is within $\cmpHeadlineAtOneFifty$ of the $401$-point value at $151$
points but $\cmpHeadlineAtHundred$ off at $101$ and $\cmpHeadlineAtFifty$ off at $51$ ($\nseq=2048$), and only
$\cmpResolvedHundred$ of its kernels are resolved at $101$ points (a registered prediction that at least $90\%$ would be
resolved failed). The safe grid is therefore about $101$ points for the $\phi^4$ kernels and about $151$ for the
mixtures, and the coarse grid cuts BP's cost by four to six times.

\begin{figure*}[t]
  \centering
  \includegraphics[width=\textwidth]{fig_grid_coarse.pdf}
  \caption{\textbf{A coarser grid.} (a)~BP time per sequence against the grid size $\ngrid$ (one thread, an
  exclusive node, $B=512$, $\nsites=32$), with the times of three networks on the same node. (b)~Relative
  change of the risk against its $401$-point value (the absolute value of the seed-mean ratio minus one) for the
  $\phi^4$ site energy and spline (base cell, $\nseq=2048$; no fit at $151$ points) and for the linear chain's
  mixture kernels with $\ncomp=1$ and $\ncomp=8$ (frozen kernels, $\nseq=2048$). Hollow markers: the coarse fit is
  \emph{better} than the $401$-point fit; points within $10^{-4}$ of it are drawn at that floor (dotted); the four series are
  offset horizontally so that coincident points can be told apart.}
  \label{fig:coarsegrid}
\end{figure*}

\paragraph{With one thread the comparison changes.}
The eight-thread ratio flatters the networks, which use threads, while BP barely does ($\cmpBpMsOneThread$\,ms on
one thread against $\cmpBpMsEightThreads$ on eight at $401$ points). On one thread and the same node a network
forward pass costs $\cmpWindowMsOne$\,ms (window), $\cmpGruMsOne$ (GRU), $\cmpUnetMsOne$ (U-Net), $\cmpTfMsOne$
(transformer) and $\cmpTcnMsOne$ (TCN), so BP at $401$ points is $\cmpRatioOneTf$ times slower than the transformer and
$\cmpRatioOneGru$ times slower than the GRU; at $151$ points ($\cmpBpMsOneFifty$\,ms) it is faster than the transformer and
TCN, within $15\%$ of the U-Net and $\cmpRatioFiftyGru$ times slower than the GRU. The networks timed are the
largest screened variants, several times the size of those scored. BP is also independent across sequences, so a
node can run it as parallel processes; that was not measured. The honest summary is that exact inference costs more than a network
forward pass by a factor that depends on the machine and the grid: $\cmpRatioCpu$ on a many-threaded CPU and
$\cmpRatioGpu$ on a GPU at $401$ points, and on one core between the cost of the largest networks and
$\cmpRatioFiftyGru$ times that of the fastest at $151$ points. What it buys with that is an error
$\matchedRatioBestD$ times smaller than the best network's at matched information on the linear chain.

\paragraph{Why not Fourier modes.}
BP's messages could be carried in Fourier space, where the convolution of the kernel with the likelihood is a
product, truncated to a few modes. We measured the premise. The message does not have a short exact description:
the natural parametric ansatz, a Gaussian passed once through the kernel, is exact at the first site only
and not a fixed point of the recursion, so its error grows into the chain by orders of magnitude (it fails to
close in all $432$ configurations tested), while it is $\cmpAnsatzOverGauss$ times closer to the true message
than a Gaussian at low noise. As a representation a Fourier basis fitted to each message needs
$\cmpFourierAdaptedLo$ numbers for a KL divergence of $10^{-2}$ against $\cmpGridLo$ for the grid, and
$\cmpFourierAdaptedHi$ for $10^{-4}$ against $\cmpGridHi$ (on the linear chain with Laplace noise), but it reaches
$10^{-4}$ in only $\cmpFourierReach$ of the cases and adding modes can make it worse, so a solver cannot pick a truncation and
know its error; the cumulant (Gram--Charlier) route is non-monotone in most cases. As a solver it would help
little: the step Fourier speeds up, the transition matrix product, is only $\cmpMatmulShareLo$--$\cmpMatmulShareHi$ of a BP
sweep at $401$ points, capping any speed-up at $\cmpFftCapLo$--$\cmpFftCapHi\times$, and the FFT transition we built is
second-order accurate ($\cmpFftError$ from dense BP) but $\cmpFftSlower$ times \emph{slower} at $401$ points, and still
$\cmpFftSlowerLarge$ times slower at the largest grid we timed ($3201$ points). Coarsening reduces every term of the sweep, which is why it is the better lever.
\Cref{app:cost} gives the tables.
